\scanpage{96}

$R\subset E$ syst. de racines ($E$ $\mathbb R$-vectoriel engendré par $R$)

produit scalaire invariant tel que $(r,r)=2$ $\forall r\in R$ (racines ttes de même longueur)

\[
\begin{aligned}
\unsure{\mathscr S|\mathscr G|\mathfrak S}&=\{x\in E^*\mid (x,r)\deleted{\tfrac{2}{(r,r)}}\in\{-1,0,1\}\ \forall r\in R\}\\
\mathscr S_0&=\{x\in E^*\mid wx-x\in R\quad\forall w\in W\ \text{t.q. }wx\ne x\}
\end{aligned}
\]

où $E^*=E-\{0\}$\\ où \deleted{$E$}$W={}$groupe de Weyl.

$\mathscr S_0\subset\mathscr S$, car $\mathscr S_0\subset E^*$, et \inserted{$x\in\mathscr S_0\Rightarrow$} $s_rx-x=-2\frac{\langle x,r\rangle}{(r,r)}r\in R$, donc si $s_rx\ne x$ i.e. $\langle x,r\rangle\ne0$, on a $-(x,r)\in\frac{(r,r)}2\{\pm1\}$, donc \deleted{\unsure{on}} tjrs \unsure{un|on|cas} $(x,r)\in\{-1,0,+1\}$.

\deleted{\unsure{Card $\mathscr S$}}

$\operatorname{Card}\mathscr S\leq3^r-1$ donc $\mathscr S$ et a fortiori $\mathscr S_0$ fini.

$\mathscr S,\mathscr S_0$ stables par $W$ et même par $\widehat W=\operatorname{Aut}(E,R)$, et $W$ a un nb fini d’orbites dans $\mathscr S$ et a fortiori dans $\mathscr S_0$.

Cas $(A_r)$

\[
\begin{aligned}
E&=\mathbb R^{r+1},\ E=\{(x_1,\ldots x_{r+1})\mid\sum_1^{r+1}x_i=0\}\\
R&=\{\alpha_{ij}=\varepsilon_i-\varepsilon_j\mid i\ne j,\ 1\leq i,j\leq r+1\}\\
\operatorname{Card} R&=r(r+1)
\end{aligned}
\]

produit scalaire induit par $(\xi,\eta)=\sum_1^{r+1}\xi_i\eta_i$ sur $\mathbb R^{r+1}$ \deleted{\unsure{est fourni par}}

\[
\begin{aligned}
&\xi_{\unsure{1|0}}\ldots\xi_r\qquad\sum\xi_i=0\\
&\underbrace{\xi_{1}=\cdots=\xi_{r_1}}_{\alpha_1}\quad
\underset{<}{\phantom{a}}\underbrace{\xi_{r_1+1}=\cdots=\xi_{r_1+r_2}}_{\alpha_2}
\quad\ldots,\quad
\underbrace{\xi_{r_1+r_2+\cdots+r_{s-1}+1}=\xi_{\underbrace{r_1+r_2+\cdots+r_s}_{r+1}}}_{\alpha_s}
\end{aligned}
\]

\[
\begin{aligned}
\left\{\begin{aligned}
r_1\alpha_1+r_2\alpha_2\ \cdots+r_s\alpha_s&=0\\
r_1+r_2\ \cdots+r_s&=r+1
\end{aligned}\right.&\qquad
\begin{aligned}\alpha_1&<\alpha_2\ \cdots<\alpha_s\\r_1,r_2&\ \cdots r_s\geq1\end{aligned}\\
\alpha_i-\alpha_j\in\{\pm1\}\quad&\forall i,j\quad\ne
\end{aligned}
\]

Soit \deleted{\illegible{2}} $s=1$, \deleted{\unsure{on}} $r_1=r+1$, $r_1\alpha_1=(r+1)\alpha_1=0$, $\alpha_1=0$ \unsure{abs.|et} donc $\xi=0$, absurde.

Soit $s=2$.

\[
\begin{aligned}
\alpha_2=\alpha_1+1\qquad&
\left\{\begin{aligned}
r_1\alpha_1+r_2\alpha_2&=0\\r_1+r_2&=r+1\\\alpha_2&=\alpha_1+1
\end{aligned}\right.
\ \Longrightarrow
\begin{aligned}
(r_1+r_2)\alpha_1+r_2&=0\\
\text{i.e. }(r+1)\alpha_1+r_2&=0\\
\text{i.e. }\alpha_1&=\frac{-r_2}{r+1}
\end{aligned}\\
\deleted{\text{Il y a }r+}\qquad&
\alpha_1=\frac{-r_2}{r+1},\ \alpha_2=\frac{r_1}{r+1}\qquad r_1+r_2=r+1
\end{aligned}
\]

il y a $r+1$ classes de tels élément, suivant la valeur de $1\leq r_1\leq r+1$

% archive sheet number: 97
